% GENERATED FILE. Edit canonical lessons, then run npm run generate:books.
% canonical-source-hash: fnv1a64:a4bf5289ce404861
% canonical-lessons: FA-C215-nomre, FA-C215-mashq, FA-C215-khabarnegar, FA-C215-bazigar, FA-C215-khanande

\chapter{Mark, Homework, Journalist, Actor, Singer}
\label{ch:fa-mark-homework-journalist-actor-singer}
\hlchaptermodality{215}

\begin{chapteropening}
\textbf{By the end of this chapter:} \emph{I can say a mark, homework, a journalist, an actor and a singer.}

Complete the last lesson of chapter 215: I can say a mark, homework, a journalist, an actor and a singer.
\end{chapteropening}

\section[\texorpdfstring{\fa{نمره} (nomre)}{نمره (nomre)}]{\fa{نمره} (nomre) — a mark, a grade}

\label{lesson:FA-C215-nomre}



\begin{quote}
Before the new one: say the Persian for a colleague, then the Persian for a manager.
\end{quote}

\subsection*{You'll want to know: \fa{نمره}}
\textbf{\fa{نمره}} — \emph{nomre} — \textquotedblleft{}a mark, a grade\textquotedblright{}.

\begin{cousinweb}
\emph{nomre} is usually linked with Italian \emph{numero}, \textquotedblleft{}number\textquotedblright{}, through Arabic or Turkish, but the route is not settled.
\end{cousinweb}

\begin{grammarlens}[title={What you've built}]
One more word for this chapter.
\end{grammarlens}

\subsection*{Your turn}
\begin{itemize}
\raggedright
  \item \emph{Say it:} \emph{nomre}
  \item \emph{Say it:} \emph{nomre}, once more
  \item \emph{Say it:} say \emph{nomre}
  \item \emph{Recall:} say the Persian for a colleague, then the Persian for a manager, then say \emph{nomre} again
\end{itemize}

\begin{tcolorbox}[breakable,colback=teal!4,colframe=teal!35!black,title={Before you move on}]
What does \fa{نمره} mean? (\textquotedblleft{}A mark, a grade\textquotedblright{}.) Say it once more.
\end{tcolorbox}
\section[\texorpdfstring{\fa{مشق} (mashq)}{مشق (mashq)}]{\fa{مشق} (mashq) — homework}

\label{lesson:FA-C215-mashq}



\begin{quote}
Before the new one: say the Persian for a manager, then the Persian for a mark.
\end{quote}

\subsection*{You'll want to know: \fa{مشق}}
\textbf{\fa{مشق}} — \emph{mashq} — \textquotedblleft{}homework\textquotedblright{}.

\begin{cousinweb}
\emph{mashq} is borrowed from Arabic \emph{mašq}, \textquotedblleft{}writing practice\textquotedblright{}. Its Arabic root is \emph{m-š-q} (Lane).
\end{cousinweb}

\begin{grammarlens}[title={What you've built}]
One more word for this chapter.
\end{grammarlens}

\subsection*{Your turn}
\begin{itemize}
\raggedright
  \item \emph{Say it:} \emph{mashq}
  \item \emph{Say it:} \emph{mashq}, once more
  \item \emph{Say it:} say \emph{mashq}
  \item \emph{Recall:} say the Persian for a manager, then the Persian for a mark, then say \emph{mashq} again
\end{itemize}

\begin{tcolorbox}[breakable,colback=teal!4,colframe=teal!35!black,title={Before you move on}]
What does \fa{مشق} mean? (\textquotedblleft{}Homework\textquotedblright{}.) Say it once more.
\end{tcolorbox}
\section[\texorpdfstring{\fa{خبرنگار} (khabarnegâr)}{خبرنگار (khabarnegâr)}]{\fa{خبرنگار} (khabarnegâr) — a journalist}

\label{lesson:FA-C215-khabarnegar}



\begin{quote}
Before the new one: say the Persian for a mark, then the Persian for homework.
\end{quote}

\subsection*{You'll want to know: \fa{خبرنگار}}
\textbf{\fa{خبرنگار}} — \emph{khabarnegâr} — \textquotedblleft{}a journalist\textquotedblright{}.

\begin{cousinweb}
\emph{khabarnegâr} is \emph{khabar}, \textquotedblleft{}news\textquotedblright{}, with \emph{-negâr}, \textquotedblleft{}writing\textquotedblright{} (from \emph{negâshtan}, \textquotedblleft{}to write, to paint\textquotedblright{}): the news-writer. \emph{khabar} is borrowed from Arabic \emph{ḫabar} (root \emph{ḫ-b-r}, Lane); \emph{negâr} is native.
\end{cousinweb}

\begin{grammarlens}[title={What you've built}]
One more word for this chapter.
\end{grammarlens}

\subsection*{Your turn}
\begin{itemize}
\raggedright
  \item \emph{Say it:} \emph{khabarnegâr}
  \item \emph{Say it:} \emph{khabarnegâr}, once more
  \item \emph{Say it:} say \emph{khabarnegâr}
  \item \emph{Recall:} say the Persian for a mark, then the Persian for homework, then say \emph{khabarnegâr} again
\end{itemize}

\begin{tcolorbox}[breakable,colback=teal!4,colframe=teal!35!black,title={Before you move on}]
What does \fa{خبرنگار} mean? (\textquotedblleft{}A journalist\textquotedblright{}.) Say it once more.
\end{tcolorbox}
\section[\texorpdfstring{\fa{بازیگر} (bâzigar)}{بازیگر (bâzigar)}]{\fa{بازیگر} (bâzigar) — an actor}

\label{lesson:FA-C215-bazigar}



\begin{quote}
Before the new one: say the Persian for homework, then the Persian for a journalist.
\end{quote}

\subsection*{You'll want to know: \fa{بازیگر}}
\textbf{\fa{بازیگر}} — \emph{bâzigar} — \textquotedblleft{}an actor\textquotedblright{}.

\begin{cousinweb}
\emph{bâzigar} is \emph{bâzi}, \textquotedblleft{}play\textquotedblright{}, with \emph{-gar}, \textquotedblleft{}doer\textquotedblright{}: the player. Both are native; \emph{bâzi} is usually traced to Middle Persian \emph{wāzīg}.
\end{cousinweb}

\begin{grammarlens}[title={What you've built}]
One more word for this chapter.
\end{grammarlens}

\subsection*{Your turn}
\begin{itemize}
\raggedright
  \item \emph{Say it:} \emph{bâzigar}
  \item \emph{Say it:} \emph{bâzigar}, once more
  \item \emph{Say it:} say \emph{bâzigar}
  \item \emph{Recall:} say the Persian for homework, then the Persian for a journalist, then say \emph{bâzigar} again
\end{itemize}

\begin{tcolorbox}[breakable,colback=teal!4,colframe=teal!35!black,title={Before you move on}]
What does \fa{بازیگر} mean? (\textquotedblleft{}An actor\textquotedblright{}.) Say it once more.
\end{tcolorbox}
\section[\texorpdfstring{\fa{خواننده} (khânande)}{خواننده (khânande)}]{\fa{خواننده} (khânande) — a singer}

\label{lesson:FA-C215-khanande}



\begin{quote}
Before the new one: say the Persian for a journalist, then the Persian for an actor.
\end{quote}

\subsection*{You'll want to know: \fa{خواننده}}
\textbf{\fa{خواننده}} — \emph{khânande} — \textquotedblleft{}a singer\textquotedblright{}.

\begin{cousinweb}
\emph{khânande} is \emph{khân-}, the present stem of \emph{khândan}, \textquotedblleft{}to read, to sing\textquotedblright{}, with \emph{-ande}, \textquotedblleft{}-ing\textquotedblright{}: the one singing. The verb is inherited.
\end{cousinweb}

\begin{grammarlens}[title={What you've built}]
One more word for this chapter.
\end{grammarlens}

\subsection*{Your turn}
\begin{itemize}
\raggedright
  \item \emph{Say it:} \emph{khânande}
  \item \emph{Say it:} \emph{khânande}, once more
  \item \emph{Say it:} say \emph{khânande}
  \item \emph{Recall:} say the Persian for a journalist, then the Persian for an actor, then say \emph{khânande} again
\end{itemize}

\begin{tcolorbox}[breakable,colback=teal!4,colframe=teal!35!black,title={Before you move on}]
What does \fa{خواننده} mean? (\textquotedblleft{}A singer\textquotedblright{}.) Say it once more.
\end{tcolorbox}
